\section{Crystalline representations and topological quantum chemistry}
\label{sec:tqc}
\subsection{Implemented inversion analysis}
An independently versioned native extension evaluates inversion
representations at time-reversal-invariant momenta (TRIM). This analysis
requires spatial inversion of the crystal, unlike the time-reversal
pairing used to reduce a scalar SCF mesh. The TRIM are diagonalized at
the converged SCF potential independently of the integration mesh,
without a Wannier90 execution or a Wannier fit. If
$T_{\mathcal I}(\bk)$ is the coefficient-space action of inversion,
including atom permutations, orbital parities and lattice phases, its
representation in the selected subspace is
\begin{equation}
 D_{\mathcal I}(\bk)=C^\dagger(\bk)S(\bk)
                   T_{\mathcal I}(\bk)C(\bk).
 \label{eq:topology-inversion}
\end{equation}
The overlap metric distinguishes the coefficient transformation
$T_{\mathcal I}$ from the covariant AO matrix of the inversion operator.
Character decomposition gives even and odd multiplicities without
assigning gauge-dependent labels to individual vectors within a mixed
degenerate eigenspace. For spinful time-reversal symmetry, the number of
odd occupied Kramers pairs yields the two-dimensional plane index or
the three-dimensional strong and weak Fu--Kane
indices~\cite{FuKane2007}, together with an inversion $z_4$ indicator
in three dimensions. An exact integer comparison with the analytic
inversion-subgroup elementary-band-representation (EBR) signatures also
tests membership in their integer lattice and the existence of a
nonnegative decomposition~\cite{Bradlyn2017,Po2017}. The implementation
uses inversion rather than a complete space-group EBR catalogue, and an
atomic-compatible signature does not establish trivial topology.
Metric, subspace, energy-commutator and time-reversal checks precede
classification. Additional parity calculations recover $\nu=0$ for
neon and $\nu=1$ for stanene. The separately versioned tests are described
in the Supporting Information and leave the retained Wilson spectra
unchanged.


\subsection{Little groups, nonsymmorphic phases, and spin}
The extension from inversion to general crystalline symmetry needs more
than a list of point rotations. Write a space-group operation as
$g=\{W_g|\tau_g\}$ in fractional direct coordinates, so that
$x\mapsto W_gx+\tau_g$. For cell matrix $A$, the Cartesian operation is
$R_g=AW_gA^{-1}$. Reciprocal fractional vectors transform with
$W_g^{-T}$, not in general with $W_g$. An antiunitary operation
$g\Theta$ additionally reverses $\bk$.

At a fixed fractional $\bk$, the little group contains operations
mapping $\bk$ to itself modulo a reciprocal vector. Choose one
representative per coset modulo integer translations of the cell.
Composition gives
\begin{equation}
 W_{gh}=W_gW_h,\qquad
 L(g,h)=\tau_g+W_g\tau_h-\tau_{gh}\in\mathbb Z^3 .
 \label{eq:coset-products}
\end{equation}
For spinors, the SU(2) lifts additionally give a sign $z(g,h)$.
The matrices acting in a fixed Bloch subspace must satisfy
\begin{equation}
 D_g\,{}^{a_g}\!D_h=\omega(g,h)D_{gh},\qquad
 \omega(g,h)=e^{-2\pi\ii\bk\cdot L(g,h)}z(g,h),
 \label{eq:projective-group}
\end{equation}
where $a_g=1$ denotes antiunitarity and
\({}^{1}\!D_h=\overline{D_h}\).
Physical $\Theta^2=-1$ contributes to the spin factor.
Associativity requires the twisted cocycle identity
\begin{equation}
 \omega(g,h)\omega(gh,j)
 ={}^{a_g}\!\omega(h,j)\,\omega(g,hj).
 \label{eq:cocycle}
\end{equation}
This distinguishes screw/glide operations from a point-group
approximation and double-valued representations from scalar ones
\cite{SpglibDefinitions,SpgrepFormulation}.

A conventional centered cell has distinct cosets with the same rotation
and different translations. Collapsing them by rotation alone loses
folded translation sectors. Conversely, character data calculated in
that conventional quotient cannot be matched blindly to a primitive-cell
EBR table. Both the cell transformation and representation convention
must be recorded.

\subsection{Metric sewing and degeneracies}
Let $T_g$ contain the atom/AO action, Bloch-image phases and spin action.
The sewing matrix between selected frames is
\begin{equation}
 B_g(\bk)=C^\dagger(g\bk)S(g\bk)
 T_g(\bk)\,{}^{a_g}\!C(\bk).
 \label{eq:sewing}
\end{equation}
For a closed selected subspace it is unitary and reconstructs the
transformed frame. Metric normalization, closure, and energy
compatibility must be checked before extracting characters.
Degeneracy is not resolved by labelling individual eigenvectors:
unitary rotations within a degenerate block are physically arbitrary.

For a unitary little group and a complete character table with the
same factor system, multiplicities follow from
\begin{equation}
 n_\alpha=\frac{1}{|G_{\bk}|}
 \sum_{g\in G_{\bk}}\overline{\chi_\alpha(g)}\,\chi(g).
 \label{eq:character-decomposition}
\end{equation}
The result must be a nonnegative integer, with consistent reconstructed
character and total dimension. Antiunitary operations require
corepresentation theory, not unmodified use of this character formula.
Characters alone also do not determine the sewing of directed Wilson
links between different k-points.

The native implementation constructs a complete projective character
table for the supplied unitary little group, rather than importing
standard-label tables. Actual GPW/GAPW Gaussian coefficient actions
include the atom permutation, AO rotation, cell-image phase and,
where requested, the post-SCF GTH spinor action. The selected band
sewing matrices are checked against Eq.~\eqref{eq:projective-group}
before characters are decomposed. Failed covariance or noninteger
multiplicities are rejected, not repaired by rotating eigenvectors
or rounding a defective representation.

\subsection{Antiunitary corepresentations}
Let the little group be $M=H\cup aH$, where $H$ is its unitary
subgroup. An irreducible projective character $\chi_\alpha$ of $H$
has antiunitary partner
\begin{equation}
 \chi_\alpha^a(h)=
 \frac{\omega(h,a)}{\omega(a,a^{-1}ha)}
 \overline{\chi_\alpha(a^{-1}ha)} .
 \label{eq:corep-partner}
\end{equation}
The Wigner indicator is~\cite{SpgrepCorepresentations}
\begin{equation}
 w_\alpha=\frac{1}{|H|}\sum_{b\in aH}
                  \omega(b,b)\chi_\alpha(b^2)\in\{+1,-1,0\}.
 \label{eq:wigner-corep}
\end{equation}
Type $+1$ retains one copy of the unitary irrep, type $-1$ requires
two copies, and type $0$ pairs two inequivalent partner irreps.
The native implementation determines these types and their integer
unitary restriction matrix. Odd type-$-1$ multiplicities or unequal
type-$0$ partner counts invalidate a selected band space.
Spin and nonsymmorphic phases enter the same factor system and cannot
be tested separately after those phases have been discarded.

These signs are neither topological indices nor antiunitary
eigenvalues. Type $+1$ does not exclude Kramers degeneracy: an
already multidimensional spin irrep can contain the required doublet
without additional doubling. At a generic centrosymmetric SOC point,
in contrast, the identity-only unitary subgroup may be doubled by
the combined inversion--time-reversal operation.
No canonical antiunitary eigenvector basis is assigned.
The electronic-structure path currently constructs the nonmagnetic
grey extension of the spatial group. The underlying algebra also
accepts supplied magnetic factor systems, but this is not a magnetic
space-group search or self-consistent noncollinear SOC calculation.

\subsection{Nonsymmorphic band compatibility on explicit segments}
For a directed connection from $\bk_a$ to $\bk_b+\bG$, define
$\mathbf d=\bk_b+\bG-\bk_a$. A segment symmetry must satisfy
\begin{equation}
 \sigma_g W_g^T\mathbf d=\mathbf d,\qquad
 \sigma_g=(-1)^{a_g},
 \label{eq:line-stabilizer}
\end{equation}
within the fractional-coordinate tolerance, not merely modulo a
reciprocal vector. Intersection of the endpoint little groups is
insufficient: ordinary time reversal can fix two TRIM endpoints
without preserving any interior point of their connecting segment.

To compare representations at the endpoints, remove the continuously
varying translational part using
\begin{equation}
 \eta_g(\bk)=e^{2\pi\ii\bk\cdot\tau_g},\qquad
 \bar\omega(g,h)=\eta_g\,{}^{a_g}\!\eta_h\,
                         \eta_{gh}^*\omega(g,h).
 \label{eq:line-rephasing}
\end{equation}
The conjugation for an antiunitary first factor is essential.
Both endpoints must produce the same rephased factor on the segment
subgroup. Its unitary projective characters provide a common basis
for integer endpoint-to-line branching matrices $B_a,B_b$.
With selected endpoint multiplicities $m_a,m_b$, compatibility requires
\begin{equation}
 B_a m_a=B_b m_b .
 \label{eq:line-compatibility}
\end{equation}
The destination rephasing uses $\bk_b+\bG$, not just $\bk_b$.
For example, a screw eigenvalue branch may exchange with another after
one reciprocal period. Ignoring $\bG$ would erase this monodromy.
Integer changes of coset representatives cancel between the
characters and the rephasing.

Endpoint corepresentations are restricted as whole representations.
If $R_a$ and $R_\ell$ give the unitary restrictions of endpoint and
line corepresentations, their branching $C_a$ satisfies
\begin{equation}
 R_\ell C_a=B_a R_a .
 \label{eq:corep-subduction}
\end{equation}
If the line has no antiunitary stabilizer, $B_aR_a$ instead describes
the branching into unitary line irreps. This retains endpoint partner
constraints even when time reversal is absent on the interior.
A spinful glide model illustrates the distinction: opposite glide
partners at Gamma become same-branch pairs at the boundary, preventing
an isolated two-band connection but allowing a four-band one.
For valid corepresentations, unitary and corepresentation line-count
compatibility are equivalent; the latter resolves the partner
structure rather than introducing a new invariant.

These operations are implemented on explicit NNKP or generated Wilson
links through \texttt{LITTLE\_GROUP\_COMPATIBILITY}, which implies
\texttt{LITTLE\_GROUP\_IRREPS}. Compatible endpoints do not prove a gap
between sampled points, a trivial Bloch bundle, or an atomic-band
decomposition. Automatic complete Brillouin-zone connectivity and
standard crystallographic irrep labels remain separate work.

\subsection{Site-induced atomic band representations}
The native induction kernel starts from a supplied fractional site $q$
and constructs its orbit $q_j=t_jq$ modulo integer translations.
Representatives $t_j$ are adjusted by lattice translations to reach the
chosen positions exactly. Localized orbitals carry irreps of the site
stabilizer; an antiunitary stabilizer instead requires the complete
corepresentations described above. Their unitary restriction characters
are denoted by $\chi_\alpha$.

For a unitary little-group operation $g$ that fixes site $j$ modulo
translations, set
\begin{equation}
 L_{gj}=W_gq_j+\tau_g-q_j\in\mathbb Z^3,\qquad
 h_j=t_j^{-1}\{I|-L_{gj}\}g t_j .
\end{equation}
The induced character is
\begin{equation}
 \chi^{\uparrow}_\alpha(g;\bk)=
 \sum_{j:\,gq_j=q_j\bmod\mathbb Z^3}
 e^{-2\pi\ii\bk\cdot L_{gj}}
 \frac{z(g,t_j)}{z(t_j,h_j)}\,
 {}^{a_{t_j}}\!\chi_\alpha(h_j).
 \label{eq:atomic-induction}
\end{equation}
This is the Bloch form of site induction~\cite{CanoBradlyn2021}, with
the spin factor in the same convention as Eq.~\eqref{eq:projective-group}.
Antiunitary orbit representatives conjugate the onsite character even
though $g$ itself is unitary. Omitting this conjugation loses complex
orbital partners.

The implementation retains conventional-cell centering cosets and keeps
each local representation column fixed across k points. Its band rank
is the orbit multiplicity times the local dimension. No full band
diagonalization or orbital gauge choice is needed to generate these
atomic reference characters. Validation compares their little-group
decompositions and segment restrictions with independent Cartesian
s/p/d orbital and spin traces. Site induction alone neither enumerates
all Wyckoff positions nor establishes that an induced representation
is elementary.

\subsection{Generating site-symmetry families}
The possible atomic reference positions need not be occupied by atoms
in the input structure. A separate geometric kernel obtains them from
the complete supplied space group. For each operation it solves
\begin{equation}
 (W_g-I)q=-\tau_g\pmod{\mathbb Z^3}
\end{equation}
and closes the resulting connected fixed sets under nonempty
intersection. Saturating the integer row lattice yields primitive
normals $B$ for each connected affine torus,
\begin{equation}
 Bq=c\pmod{\mathbb Z^r},\qquad d=3-r.
\end{equation}
Distinct components are retained: for example, $2q_x=0\pmod1$
contains both $q_x=0$ and $q_x=1/2$. Integer right-hand sides are
enumerated from the bounds of the unit cube rather than by sampling
a coordinate mesh. Subsequent space-group equivalence identifies the
families, and a generic representative is checked against its entire
site stabilizer.

Closure containment supplies specialization relations between families.
A site group is maximal when its family has no higher-symmetry
specialization. This provides the geometric input to site induction,
but does not replace the representation-level elementarity tests needed
for an EBR catalogue. The construction retains the supplied cell and
origin conventions and does not assign standard Wyckoff letters.

\subsection{Atomic signature matrices on explicit reciprocal sets}
Combining the generated sites with the induction kernel gives a native
integer signature matrix for a supplied reciprocal set $\mathcal K$,
\begin{equation}
 A_{(\bk,\beta),(q,\alpha)}
 =n_{\beta}\!\left[\chi^{\uparrow}_{q,\alpha}(\cdot;\bk)\right],
 \qquad \bk\in\mathcal K .
\end{equation}
The column $(q,\alpha)$ retains one Wyckoff-family representative and
one onsite irrep or corepresentation at every point. Rows use unitary
irreps unless the little group has an antiunitary coset; in that case
they count complete corepresentations after checking the required
unitary partner multiplicities. Consequently,
\begin{equation}
 \sum_\beta d_{\bk\beta} A_{(\bk,\beta),(q,\alpha)}
 =N_{q\alpha}
\end{equation}
is the same physical band rank for every $\bk$. Counts are not generally
numbers of Kramers pairs: corepresentation dimensions and pairing types
are retained explicitly.

The implementation preserves all sites, including nonmaximal ones,
and does not identify columns merely because their sampled signatures
coincide. An incoming unitary character vector can be decomposed in
the same row basis after matching its operation indices, provided it
uses the same spin/Bloch convention. This algebraic check does not replace
Hamiltonian covariance or selected-subspace closure for actual Gaussian
states. The matrix remains a redundant atomic reference for the
explicit set $\mathcal K$, not an automatically complete EBR matrix:
reciprocal-stratum completeness and representation elementarity remain
separate from the integer decomposition of this supplied matrix.

\subsection{Exact signed and nonnegative signature decompositions}
For an integer signature matrix $A\in\mathbb Z^{m\times n}$, the native
integer kernel constructs a Smith form~\cite{CanoBradlyn2021},
\begin{equation}
 UAV=D,\qquad
 D=\operatorname{diag}(d_1,\ldots,d_r,0,\ldots),\qquad
 d_i>0,\quad d_i\mid d_{i+1},
\end{equation}
where $U$ and $V$ are unimodular integer matrices. With $c=Ub$,
the equation $Ax=b$ has an integer solution precisely when
\begin{equation}
 c_i=0\quad(i>r),\qquad c_i\equiv0\pmod{d_i}\quad(i\leq r).
\end{equation}
A witness is obtained as $x=Vy$, with $y_i=c_i/d_i$ for $i\leq r$
and zero free components. The implementation checks both $UAV=D$
and $Ax=b$ using bounded 64-bit integer arithmetic. Rank decisions
do not use floating-point thresholds. Growth beyond the arithmetic
bound or exhaustion of the factorization work limit returns an
unresolved status, not a negative membership decision.

For the nonnegative signature matrices used here, the stronger condition
$Ax=b$, $x\in\mathbb Z_{\geq0}^n$, is tested separately. At a search
node with residual $b'\geq0$, a nonzero column $j$ has the finite bound
\begin{equation}
 0\leq x_j\leq
 \min_{i:A_{ij}>0}\left\lfloor\frac{b'_i}{A_{ij}}\right\rfloor.
\end{equation}
Row-wise gcds of the remaining columns and their zero supports prune
the complete enumeration; zero columns can be set to zero. Every
accepted witness is checked exactly. A caller-specified node budget
bounds the potentially exponential search. An exhausted budget is
explicitly undecided and must not be interpreted as absence of an
atomic decomposition.

These are decisions about the supplied columns and sampled rows.
The Smith divisors of an incomplete reciprocal sample are not, by
themselves, the space group's symmetry-indicator group. Nonnegative
membership likewise does not prove triviality of a Bloch bundle.
The integer routines do not add an automatic general-classification
input mode or replace spectral isolation and compatibility checks.

\subsection{Band representations and limits of indicators}
Elementary band representations (EBRs) provide the indecomposable
building blocks of atomic band representations. Restricting them to
little groups yields a vector of irrep multiplicities. Compatibility
relations constrain the vectors that can connect along high-symmetry
lines. A candidate symmetry vector $n$ can be compared with an integer
matrix of EBR signatures,
\begin{equation}
 E q=n .
 \label{eq:ebr-system}
\end{equation}
Integer solvability, nonnegative solvability, and compatibility
constraints are different questions. A symmetry indicator is a
quotient of compatible representation data by atomic data; it can
miss topology invisible to those symmetry eigenvalues
\cite{Bradlyn2017,Po2017}.
An atomic-compatible signature is not proof of a trivial Bloch bundle.
A signed-only decomposition is not, without additional information,
a complete diagnosis of fragile topology.

The implemented EBR-signature classifier still uses the inversion
subgroup and its analytic atomic signatures. The general native
little-group/corepresentation, explicit-segment, Wyckoff-family and site-induction
machinery supplies additional representation data, not a complete
all-space-group classification. Conventional-to-primitive matching,
standard labels, complete connectivity and a validated catalogue with
elementarity checks are still required. Algebra tests over all ordinary
space groups must not be presented as a completed classification
of all their materials.

\section{Parameter-space topology and phasons}
\label{sec:phasons}
A separate native postprocessor, \texttt{topology\_phasons}, extends the
overlap construction to closed two- or four-parameter families, including
periodic atomic displacements (phasons) combined with reciprocal
coordinates. It reads \texttt{STATE\_EXPORT} snapshots and evaluates
$M_{ab}=C_a^\dagger O_{ab}C_b$, where $O_{ab}$ includes the Gaussian
centers of both geometries, periodic images, and Bloch and boundary-sewing
phases. Polar-unitary links give $C_1$ from determinant plaquette
phases~\cite{Fukui2005} and a four-parameter $C_2$ estimate from full
non-Abelian plaquette logarithms~\cite{MocholGrzelak2019}. Here $C_2$
denotes the second Chern character pairing, not generally the second
Chern class. The postprocessor requires neither Python, Z2Pack, nor
Wannier90 at runtime. It operates on fixed-cell, fixed-rank families with
explicit physical endpoint identifications. The Supporting Information
gives its separate development version, discretization, and model and
Gaussian-snapshot tests. These tests complement the material examples
below without changing their Hamiltonians, sampling, or results.


For a discrete parameter family, let $Q_\mu(x)$ denote the polar link
from $x$ to $x+\hat\mu$. An oriented plaquette is
\begin{equation}
 P_{\mu\nu}(x)=Q_\mu(x)Q_\nu(x+\hat\mu)
 Q_\mu^\dagger(x+\hat\nu)Q_\nu^\dagger(x).
\end{equation}
With the implemented directed-overlap convention,
$C_1^{\mathrm{mesh}}=-(2\pi)^{-1}\sum_x\arg\det P_{12}(x)$.
For four parameters write $F_{\mu\nu}(x)=\log P_{\mu\nu}(x)$
using the principal matrix logarithm. The non-Abelian estimator is
\begin{equation}
 C_2^{\mathrm{mesh}}=-\frac{1}{4\pi^2}\sum_x
 \Re\Tr\!\left[F_{12}F_{34}-F_{13}F_{24}+F_{14}F_{23}\right](x).
 \label{eq:phason-c2}
\end{equation}
The logarithms already include the discrete plaquette area.
Mesh refinement, physical endpoint identifications, and a fixed isolated
rank are essential. Coefficient arrays from different geometries
cannot be contracted without cross-geometry AO overlaps.
The numerical value is retained before rounding, and a sampled gap
does not certify isolation between sampled geometries.
